En este capitulo se resume la estrategia de validacion utilizada en la iteracion actual del proyecto. Dado que se trata de un \textit{TFG} con una aplicacion \textit{full-stack} de alcance incremental, la verificacion combina comprobaciones automaticas de regresion con validaciones funcionales manuales y asistidas sobre los flujos criticos de los modulos \texttt{M1}, \texttt{M2}, \texttt{M3} y \texttt{M4}.

\section{Estrategia}

La estrategia de pruebas seguida persigue dos objetivos complementarios:

\begin{itemize}
    \item detectar regresiones tecnicas rapidas despues de cada cambio relevante mediante comprobaciones automaticas;
    \item confirmar de forma funcional que los recorridos visibles para cada rol continuan funcionando de acuerdo con el alcance operativo implementado.
\end{itemize}

Las regresiones tecnicas se evaluan repitiendo siempre la misma secuencia base:

\begin{enumerate}
    \item comprobacion estatica de tipos con \texttt{TypeScript};
    \item analisis estatico de codigo con \texttt{ESLint};
    \item construccion completa de produccion con \texttt{Next.js};
    \item smoke test funcional de \texttt{M4} sobre servicios y persistencia;
    \item pasada visual hidratada de \texttt{M4} sobre rutas criticas por rol.
\end{enumerate}

Esta rutina ofrece una senal rapida sobre errores de tipado, problemas de compilacion, importaciones rotas, uso incorrecto de componentes de servidor o cliente, fallos en rutas dinamicas y desviaciones funcionales del flujo de pedidos, reparto y cobro.

\section{Pruebas Unitarias}

En la version actual no se ha incorporado todavia una bateria formal de pruebas unitarias con un \textit{runner} dedicado. Aun asi, parte de la validacion de bajo nivel se apoya en tres mecanismos que reducen errores locales antes de integrar cambios:

\begin{itemize}
    \item el tipado estricto de \texttt{TypeScript} sobre contratos, componentes y servicios;
    \item la centralizacion de reglas de negocio en servicios de dominio de usuarios, operativa comercial, catalogo y pedidos;
    \item las validaciones normalizadas de entrada antes de persistir cambios, lo que limita estados inconsistentes en operaciones de administracion, catalogo, reparto y cobro.
\end{itemize}

La ausencia de una suite unitaria completa sigue siendo una mejora pendiente, pero la separacion actual de servicios facilita su incorporacion futura.

\section{Pruebas de Integracion}

Las pruebas de integracion se centran en comprobar la interaccion entre interfaz, \textit{Route Handlers}, autenticacion, persistencia y servicios externos. En esta linea resultan especialmente relevantes los siguientes puntos:

\begin{itemize}
    \item autenticacion con \texttt{Auth.js} y control de acceso por rol en \texttt{proxy.ts};
    \item coordinacion entre formularios administrativos y servicios \texttt{TypeORM} para usuarios, clientes, visitas, catalogo y pedidos;
    \item integracion con \texttt{Nominatim} y \texttt{OSRM} para geocodificacion y calculo de rutas del modulo \texttt{M2};
    \item integracion con \texttt{Cloudinary} para subida, reemplazo y validacion de imagenes de perfil y catalogo en \texttt{M1} y \texttt{M3};
    \item integracion con \texttt{QuickChart} para representar el QR operativo de cada pedido dentro del flujo de reparto de \texttt{M4}.
\end{itemize}

Estas comprobaciones se realizan principalmente a traves de la propia aplicacion en entorno local, verificando que una accion iniciada desde la interfaz produce el efecto esperado en la API y en la base de datos.

\section{Pruebas de Sistema}

La evidencia automatizada disponible para el estado actual del repositorio se ha obtenido ejecutando los siguientes comandos sobre la aplicacion \texttt{app-tfg/}:

\begin{itemize}
    \item \texttt{npm run typecheck}
    \item \texttt{npm run lint}
    \item \texttt{npm run build}
    \item \texttt{npm run m4:closeout}
    \item \texttt{npm run m4:ui-check}
\end{itemize}

El resultado de esta validacion es satisfactorio: la comprobacion de tipos, el analisis estatico y la construccion completa de produccion finalizan sin errores, y el cierre tecnico de \texttt{M4} queda reforzado por dos evidencias adicionales:

\begin{itemize}
    \item un smoke test funcional de servicios y persistencia, que cubre el flujo principal de pedidos, reparto y cobro y finaliza con \texttt{19/19} comprobaciones satisfactorias en la bateria principal y \texttt{12/12} en la pasada de cierre;
    \item una verificacion visual hidratada sobre cliente, comercial y administrador, reutilizable para comprobar bandeja de reparto, detalle de visita, cobros, ETA del cliente, detalle administrativo y gestion de errores de sesion o de carga.
\end{itemize}

Una mejora tecnica asociada a esta fase de pruebas ha consistido en desacoplar la construccion de \texttt{Next.js} de fuentes remotas, de modo que \texttt{npm run build} ya no depende de descargar \texttt{Geist} desde \texttt{next/font/google} para completarse en local.

\subsection{Pruebas Funcionales}

Las comprobaciones funcionales prioritarias se organizan por modulo:

\begin{itemize}
    \item \texttt{M1}: registro de solicitud de alta, autenticacion con credenciales, navegacion protegida por rol, edicion de perfil y cambio de contrasena.
    \item \texttt{M2}: administracion de clientes, asignacion comercial-cliente, geolocalizacion confirmada, creacion y seguimiento de visitas, configuracion diaria del comercial y calculo de ETA para cliente.
    \item \texttt{M3}: mantenimiento del catalogo desde administracion, consulta publica de productos activos, filtrado por categorias y lineas, acceso a fichas tecnicas, consulta de recursos de apoyo y exploracion de cartas de color con sus referencias.
    \item \texttt{M4}: borradores de pedido en cliente y comercial, confirmacion y cancelacion coherente, bandeja de pedidos pendientes de reparto, visita de entrega con validacion QR, paso a \texttt{delivered}, subflujo minimo de cobro y trazabilidad cruzada desde admin.
\end{itemize}

En \texttt{M4}, la regresion funcional mas sensible consiste en comprobar que el mismo pedido se mantiene coherente al pasar por tres contextos distintos: cliente, comercial y administrador. Por ello, cualquier cambio en entidades, serializadores, servicios de pedido o logica de visitas de reparto debe revalidarse en esos tres recorridos.

\subsection{Pruebas No Funcionales}

En el estado actual del proyecto, las comprobaciones no funcionales se concentran en:

\begin{itemize}
    \item seguridad basica de acceso mediante autenticacion y autorizacion por roles;
    \item robustez de compilacion y tipado de toda la aplicacion;
    \item validacion basica de errores de sesion y de carga en rutas criticas de \texttt{M4};
    \item compatibilidad minima de navegador a traves de la ruta informativa para clientes no soportados;
    \item validacion de tamano, tipo y procedencia de imagenes subidas a \texttt{Cloudinary}.
\end{itemize}

No se ha abordado todavia una campana especifica de rendimiento, carga o pruebas automatizadas de interfaz con un \textit{framework} dedicado, por lo que estas areas permanecen como trabajo de mejora futura.

\section{Pruebas de Aceptacion}

La aceptacion de la version actual del sistema se interpreta como la disponibilidad integrada y estable de los modulos \texttt{M1}, \texttt{M2}, \texttt{M3} y \texttt{M4}. En terminos practicos, esto implica que:

\begin{itemize}
    \item el administrador puede gestionar usuarios, clientes, catalogo y pedidos globales;
    \item el comercial puede consultar su operativa diaria, el catalogo, sus pedidos, los repartos asociados y el subflujo minimo de cobro;
    \item el cliente puede consultar su prevision de reparto, su perfil, el catalogo, la coloracion y el historial de pedidos;
    \item el flujo \texttt{pedido -> reparto -> entrega -> cobro minimo} queda cerrado y trazable en la implementacion actual.
\end{itemize}

Con este umbral cubierto, el siguiente salto funcional razonable del proyecto no pasa ya por incorporar \texttt{M4}, sino por consolidar documentalmente su alcance final y, en futuras iteraciones, ampliar el dominio de cobros hacia pagos parciales, facturacion y conciliacion avanzada si el proyecto lo requiere.